\section{Introduction }

This chapter describes the main threats to validity and limitations of the AI-Scientist thesis. The goal is not to
weaken the contribution, but to state precisely where the evidence applies. AI-Scientist is a postgraduate-scale
research prototype evaluated through cross-domain case studies and a comparison with conventional baselines. Its
strongest claims concern evidence grounding, calibrated confidence, contradiction handling, measured abstention,
and traceability. Its broader applicability depends on the quality of the connected evidence sources and the scope
of the retrieval design.

The limitations are organized using four standard categories: construct validity, internal validity, external
validity, and conclusion validity. The chapter then summarizes the practical engineering limitations of the current
system.

\section{Construct Validity }

Construct validity asks whether the thesis measures the concepts it claims to measure. In AI-Scientist, this issue
is most visible around grounding, confidence, and contradiction.

\subsection{Grounding Is Not Correctness }

The thesis treats a verdict as trustworthy when it is bound to retrieved evidence. However, a grounded verdict is
not equivalent to scientific truth. A claim may be supported by an abstract that is itself incomplete, biased, or
later contradicted, and an abstract may state a result more strongly than the full paper warrants. Grounding
therefore certifies consistency with the retrieved evidence while remaining distinct from full scientific
validation, which is a standard distinction in retrieval-grounded evaluation \cite{petroni2020kilt,thakur2021beir,es2023ragas}.

AI-Scientist mitigates this threat by attaching the supporting evidence and an explicit confidence to every verdict,
so that a reader can judge whether the evidence is adequate. Even so, this safeguard is intentionally scoped: it
makes the basis for a conclusion visible, while leaving room for a reader to assess the underlying evidence base.

\subsection{Abstracts Are a Proxy for Papers }

The system reasons primarily over titles, abstracts, and uploaded paper records rather than full text. Abstracts are
a reasonable proxy because they state a paper's main claims, and they keep the system efficient and broadly
applicable. But they are a proxy nonetheless. A nuance, caveat, or limitation stated only in the body of a paper
will not be visible to the verifier, which can make a claim look better supported than the full text would justify.
This proxy-based setup is common in scientific QA and factuality benchmarks, where concise evidence is often the
practical starting point \cite{thorne2018fever,wadden2020scifact,huang2023hallucinationsurvey}.

The thesis is explicit that full-text deep reading is a non-goal, so this is a scoping decision rather than a hidden
issue. Nevertheless, conclusions should be read as grounded in abstract-level evidence, complemented by the source
metadata available in the system, which mirrors the abstract-centric evidence setup used in several retrieval and
fact-verification benchmarks \cite{thorne2018fever,wadden2020scifact,petroni2020kilt}.

\subsection{Token Overlap Is a Proxy for Semantic Support }

Evidence matching uses token-overlap and coverage measures rather than deep semantic entailment. These measures are
transparent, reproducible, and cheap, and the morphological normalization helps related word forms match. However,
lexical overlap can both miss genuine support phrased in different words and over-credit superficial term matches
that do not actually entail the claim. The contradiction analysis, which relies on negation and cue detection, has
the same character: it captures common polarity mismatches but cannot detect every subtle disagreement.
That trade-off is typical in retrieval-first systems, which often use lexical signals as a reliable first pass before
more expensive semantic reasoning \cite{karpukhin2020dpr,khattab2020colbert,guu2020realm}.

The optional language-model refinement partially addresses this by allowing a model to adjust borderline verdicts
over the already-retrieved evidence, but the deterministic core remains lexical. Verdicts should therefore be
understood as evidence-overlap judgments, strengthened by optional model refinement, rather than as full semantic
entailment decisions, a common trade-off in retrieval-first QA systems \cite{karpukhin2020dpr,khattab2020colbert,guu2020realm}.

\subsection{Confidence Is a Designed Heuristic }

Confidence is computed as a transparent function of evidence strength, corroboration, and source independence,
clamped between a floor and a ceiling. This makes it inspectable and monotonic in the intended factors, but the
specific coefficients are design choices, not learned calibration against labelled outcomes. The confidence is best
read as a relative, well-behaved ordering --- higher means better corroborated --- rather than as a probability with
a guaranteed frequentist interpretation.

\section{Internal Validity }

Internal validity asks whether the observed behaviour is caused by the intended design factors rather than by
incidental artifacts.

\subsection{Shared Substrate for a Fair Comparison }

The baselines are implemented on the same retrieval and scoring substrate as the full system, so the comparison
isolates orchestration, claim-level verification, and critique. This is a deliberate control. The residual risk is
that the shared substrate itself shapes all three systems, so the comparison should be read as measuring the value
of the multi-agent layer \emph{given} this retrieval foundation, not as an absolute statement about every possible
retriever. This is consistent with benchmark practice in retrieval-oriented evaluation settings
\cite{petroni2020kilt,thakur2021beir,ragsurvey2024,min2023factscore}.

\subsection{Model Non-Determinism }

When the optional language model is enabled, its outputs are not perfectly reproducible. Provider-side updates,
decoding variation, and transient failures can change a refined verdict or a rewritten summary. AI-Scientist
mitigates this by keeping the grounding, verdict, and confidence logic deterministic and by treating the model as an
optional refinement that falls back to the deterministic result on failure. The deterministic behaviour is therefore
reproducible; the model-refined behaviour carries the usual variability of hosted models.

\subsection{Live-Source Variability }

Live retrieval introduces variability that is inherent rather than accidental. The same question asked on different
days may surface a slightly different set of current papers, because PubMed Central and arXiv are themselves
changing. This is disclosed rather than hidden. For reproducible inspection, the curated corpus and uploaded papers
provide a stable evidence base, while live retrieval is the component whose results legitimately vary over time.

\section{External Validity }

External validity asks whether the behaviour generalizes beyond the evaluated questions, domains, and sources.

\subsection{Domain and Source Coverage }

The evaluation spans medicine, computer science, biology, and psychology, served by PubMed Central and arXiv. These
are broad and important domains, but they do not cover every field, and the two live sources do not index every
discipline equally. Questions in areas poorly served by these sources --- for example, parts of the humanities, law,
or proprietary engineering literature --- would retrieve weaker evidence, and the system's grounding is only as good
as the sources it can reach. That limitation is typical of live retrieval systems, where coverage depends on source
availability and indexing depth \cite{gu2021pubmedbert,lee2019biobert,ragsurvey2024,petroni2020kilt}.

The architectural ideas --- layered evidence, claim-level verification, calibrated confidence, and critique --- are
source-agnostic in principle. The measured effectiveness, however, depends on the coverage and quality of the
connected sources.

\subsection{Question Phrasing and Claim Selection }

The system extracts claims from retrieved abstracts using surface cues and question-term overlap. This works well
for questions phrased in the vocabulary of the field, but an unusually phrased question, or one using terminology
that differs from the abstracts, may yield weaker extraction and retrieval. The behaviour is therefore somewhat
sensitive to phrasing, and the case studies, while cross-domain, cannot represent every possible way a question
could be asked. Similar phrasing sensitivity is a known challenge in retrieval and fact-verification pipelines
\cite{min2023factscore,lin2021truthfulqa,honovich2022true,es2023ragas}.

\subsection{Reliance on Live Availability }

Part of the system's value comes from current live evidence. In settings without network access, or where a source
is rate-limited or unavailable, the system falls back to the curated corpus and uploads. This fallback keeps it
usable, but it narrows the evidence base, so the external validity of the live-augmented behaviour is contingent on
the sources being reachable. The thesis therefore treats live retrieval as an enhancement to a stable grounded core
rather than as a requirement for correctness \cite{lewis2020rag,guu2020realm,asai2024selfrag}.

\section{Conclusion Validity }

Conclusion validity asks whether the evidence is strong enough to support the stated conclusions.

\subsection{Case-Study Evaluation }

The evaluation is qualitative and case-study based rather than a large aggregate score over a labelled benchmark.
This is appropriate for a verification system whose claims are about the structure and clarity of its reasoning,
which a single accuracy number would obscure. The trade-off is that the conclusions are demonstrations of behaviour
on representative questions rather than statistical generalizations. The thesis is careful to phrase its claims
accordingly: it argues that the system grounds, calibrates, and abstains as designed, illustrated across domains,
rather than asserting a precise success rate, matching the evaluation style used in recent grounded QA research
\cite{es2023ragas,min2023factscore,lin2021truthfulqa}.

\subsection{Designer-Set Thresholds }

Several behaviours depend on thresholds: the minimum coverage for retrieval, the support and contradiction overlap
levels, the confidence floor and ceiling, and the critic's confidence threshold. Thresholds are necessary for a
deployable system, but the exact defaults are design choices and could be tuned. The thesis treats them as
transparent, inspectable operating points rather than as universal constants, and their values are centralized in
the configuration so that their influence is visible.

\section{Practical System Limitations }

AI-Scientist is not a production research engine in the broad sense. It demonstrates an end-to-end architecture and
the qualities claimed, but several engineering limitations remain.

First, the system reasons over abstracts and uploaded records rather than full papers, so it cannot assess
methodological details, sample sizes, or statistical rigour that appear only in the body of a paper.

Second, evidence matching is primarily lexical, so it can miss paraphrased support and subtle contradictions that a
deeper semantic model would catch.

Third, live coverage is limited to the connected sources, and the quality of grounding varies with how well a domain
is indexed by PubMed Central and arXiv.

Fourth, the confidence and verdict thresholds are designer-set heuristics rather than learned, calibrated
parameters, so confidence should be read as a well-behaved ordering rather than a probability.

\begin{table}[!htbp]
\centering
\caption{Summary of the principal limitations and their effect}
\begin{tabularx}{0.95\linewidth}{>{\raggedright\arraybackslash}p{0.26\linewidth} >{\raggedright\arraybackslash}X >{\raggedright\arraybackslash}p{0.22\linewidth}}
\toprule
Limitation & Effect on the system & Practical implication \\
\midrule
Abstract-level evidence & Full-paper nuance may be hidden & Results should be read at the abstract level \\
Lexical matching & Paraphrased support can be missed & Better semantic matching would strengthen coverage \\
Live-source dependence & Coverage varies by source availability & Offline corpus remains an important fallback \\
Heuristic thresholds & Confidence is not probabilistic & Scores should be interpreted comparatively \\
\bottomrule
\end{tabularx}
\end{table}

\begin{figure}[!htbp]
\centering
\begin{tikzpicture}[
  font=\small,
  scale=0.9,
  transform shape,
  box/.style={draw, rounded corners, align=center, text width=28mm, minimum height=10mm, fill=blue!8},
  arrow/.style={-{Stealth[round]}, thick}
]
\node[box] (a) {Grounded prototype};
\node[box, right=6mm of a] (b) {Evidence limits};
\node[box, right=6mm of b] (c) {Reader interpretation};
\node[box, right=6mm of c] (d) {Future refinement};

\draw[arrow] (a) -- (b);
\draw[arrow] (b) -- (c);
\draw[arrow] (c) -- (d);
\end{tikzpicture}
\caption{How the current limitations guide future refinement}
\label{fig:limitations_flow}
\end{figure}

\section{Chapter Conclusion }

The main limitations of AI-Scientist are its reliance on abstracts rather than full text, its lexical evidence
matching, its dependence on live-source coverage, and the case-study nature of its evaluation. These limitations do
not invalidate the central contribution, because the thesis is careful about what it claims: AI-Scientist is an
evidence-grounded, self-critical verification assistant that reports its uncertainty clearly and keeps its
conclusions aligned with the evidence it retrieves. The next chapter concludes the thesis and outlines future work
that follows directly from these limits.

